% attack17: referee round 5, comment 1 -- rigorous membership proof.
% Paste-ready LaTeX paragraph.  Assumes \cite{Brunault} = arXiv:1504.08127
% (J. Number Theory 163 (2016)) and that F_total has been defined as the
% finite signed sum over the seven Manin symbols of the products
% e_{a,d}e_{b,-c}+e_{a,-d}e_{b,c}.  All computational claims refer to
% code/siegel_anchor_step12.py and code/siegel_anchor_step12_out.json.

\paragraph{Membership of $F_{\mathrm{total}}$ and the Sturm argument.}
Each series $e_{a,b}$, $(a,b)\in(\mathbb{Z}/11\mathbb{Z})^2$, occurring in the
definition of $F_{\mathrm{total}}$ is, by \cite[Definition~10 and
Lemma~11]{Brunault}, an Eisenstein series of weight~$1$ on
$\Gamma_1(11^2)=\Gamma_1(121)$; in particular it is a modular form, holomorphic
on $\mathcal{H}$ and at every cusp.  Every product
$e_{a,d}e_{b,-c}+e_{a,-d}e_{b,c}$ therefore lies in $M_2(\Gamma_1(121))$, and
since $F_{\mathrm{total}}$ is by construction a finite signed sum of such
products (seven Manin symbols, each contributing finitely many terms), we have
the \emph{unconditional} membership
\[
  F_{\mathrm{total}} \;\in\; M_2\bigl(\Gamma_1(121)\bigr).
\]
As $-2f_{11}\in M_2(\Gamma_0(11))\subset M_2(\Gamma_1(121))$, the difference
$D:=F_{\mathrm{total}}+2f_{11}$ also lies in $M_2(\Gamma_1(121))$.  The Sturm
bound for this space is
\[
  \frac{k}{12}\,\bigl[\mathrm{PSL}_2(\mathbb{Z}):\bar\Gamma_1(121)\bigr]
  \;=\; \frac{2}{12}\cdot 7260 \;=\; 1210,
  \qquad
  \bigl[\mathrm{PSL}_2(\mathbb{Z}):\bar\Gamma_1(121)\bigr]
  = \frac{121^2}{2}\Bigl(1-\frac{1}{121}\Bigr) = 7260,
\]
where $-I$ acts trivially because the weight is even.  We computed the
$q$-expansion of $D$ in exact rational arithmetic (the coefficients of
$e_{a,b}$ at $q^n$, $n\ge1$, are integers and $\alpha_0(a,b)\in\frac1{22}\mathbb{Z}$,
so all products have rational coefficients; the convolutions were evaluated
exactly), and found
\[
  a_n(D)=0 \qquad\text{for all } 0\le n\le 2420,
\]
twice the sharp Sturm bound and also exceeding the conservative
$\mathrm{SL}_2$-index convention $\frac{2}{12}\cdot 14520 = 2420$
(\texttt{code/siegel\_anchor\_step12.py}, output
\texttt{siegel\_anchor\_step12\_out.json}: PASS; cross-validated against the
$251$ exact coefficients of the earlier computation).  By Sturm's theorem
$D=0$, that is,
\[
  F_{\mathrm{total}} \;=\; -2\,f_{11}
\]
as modular forms.  A fortiori $F_{\mathrm{total}}\in M_2(\Gamma_0(11))$, so the
subsequent use of the Sturm bound $2$ for $M_2(\Gamma_0(11))$ is legitimate and
the argument is not circular.

\begin{remark}[Per-symbol decomposition]
The same Sturm computation applied to each of the seven Manin symbols
individually (each is a finite sum of products of the $e_{a,b}$, hence lies in
$M_2(\Gamma_1(121))$) certifies that the symbols $\{0,3/11\}$-component
$[[1,0],[3,1]]$ and $[[1,2],[1,3]]$ each contribute exactly $-f_{11}$, while
the other five symbols are identically zero as modular forms (coefficients
checked exactly through $q^{2420}$).  The main proof above uses only the total
$F_{\mathrm{total}}=-2f_{11}$; the per-symbol decomposition is supplementary.
\end{remark}
